% Options for packages loaded elsewhere
%DIF LATEXDIFF DIFFERENCE FILE
%DIF DEL D:\research\quantum_paper\paper\mdpi_sensors\revision\response_to_editor_final.tex           Sun Aug 16 17:29:56 2026
%DIF ADD D:\research\quantum_paper\paper\mdpi_sensors\revision\response_to_editor_removal_final.tex   Sun Aug 16 18:27:23 2026
\PassOptionsToPackage{unicode}{hyperref}
\PassOptionsToPackage{hyphens}{url}
\documentclass[
  11pt,
]{article}
\usepackage{xcolor}
\usepackage[margin=1in]{geometry}
\usepackage{amsmath,amssymb}
\setcounter{secnumdepth}{-\maxdimen} % remove section numbering
\usepackage{iftex}
\ifPDFTeX
  \usepackage[T1]{fontenc}
  \usepackage[utf8]{inputenc}
  \usepackage{textcomp} % provide euro and other symbols
\else % if luatex or xetex
  \usepackage{unicode-math} % this also loads fontspec
  \defaultfontfeatures{Scale=MatchLowercase}
  \defaultfontfeatures[\rmfamily]{Ligatures=TeX,Scale=1}
\fi
\usepackage{lmodern}
\ifPDFTeX\else
  % xetex/luatex font selection
  \setmainfont[]{Times New Roman}
\fi
% Use upquote if available, for straight quotes in verbatim environments
\IfFileExists{upquote.sty}{\usepackage{upquote}}{}
\IfFileExists{microtype.sty}{% use microtype if available
  \usepackage[]{microtype}
  \UseMicrotypeSet[protrusion]{basicmath} % disable protrusion for tt fonts
}{}
\makeatletter
\@ifundefined{KOMAClassName}{% if non-KOMA class
  \IfFileExists{parskip.sty}{%
    \usepackage{parskip}
  }{% else
    \setlength{\parindent}{0pt}
    \setlength{\parskip}{6pt plus 2pt minus 1pt}}
}{% if KOMA class
  \KOMAoptions{parskip=half}}
\makeatother
\setlength{\emergencystretch}{3em} % prevent overfull lines
\providecommand{\tightlist}{%
  \setlength{\itemsep}{0pt}\setlength{\parskip}{0pt}}
\usepackage{bookmark}
\IfFileExists{xurl.sty}{\usepackage{xurl}}{} % add URL line breaks if available
\urlstyle{same}
\hypersetup{
  hidelinks,
  pdfcreator={LaTeX via pandoc}}

\author{}
\date{}

%DIF PREAMBLE EXTENSION ADDED BY LATEXDIFF
%DIF UNDERLINE PREAMBLE %DIF PREAMBLE
\RequirePackage[normalem]{ulem} %DIF PREAMBLE
\RequirePackage{color}\definecolor{RED}{rgb}{1,0,0}\definecolor{BLUE}{rgb}{0,0,1} %DIF PREAMBLE
\providecommand{\DIFadd}[1]{\hl{#1}} %DIF PREAMBLE
\providecommand{\DIFdel}[1]{} %DIF PREAMBLE
%DIF SAFE PREAMBLE %DIF PREAMBLE
\providecommand{\DIFaddbegin}{} %DIF PREAMBLE
\providecommand{\DIFaddend}{} %DIF PREAMBLE
\providecommand{\DIFdelbegin}{} %DIF PREAMBLE
\providecommand{\DIFdelend}{} %DIF PREAMBLE
\providecommand{\DIFmodbegin}{} %DIF PREAMBLE
\providecommand{\DIFmodend}{} %DIF PREAMBLE
%DIF FLOATSAFE PREAMBLE %DIF PREAMBLE
\providecommand{\DIFaddFL}[1]{\DIFadd{#1}} %DIF PREAMBLE
\providecommand{\DIFdelFL}[1]{\DIFdel{#1}} %DIF PREAMBLE
\providecommand{\DIFaddbeginFL}{} %DIF PREAMBLE
\providecommand{\DIFaddendFL}{} %DIF PREAMBLE
\providecommand{\DIFdelbeginFL}{} %DIF PREAMBLE
\providecommand{\DIFdelendFL}{} %DIF PREAMBLE
%DIF AMSMATHULEM PREAMBLE %DIF PREAMBLE
\makeatletter %DIF PREAMBLE
\let\sout@orig\sout %DIF PREAMBLE
\renewcommand{\sout}[1]{\ifmmode\text{\sout@orig{\ensuremath{#1}}}\else\sout@orig{#1}\fi} %DIF PREAMBLE
\makeatother %DIF PREAMBLE
%DIF COLORLISTINGS PREAMBLE %DIF PREAMBLE
\RequirePackage{listings} %DIF PREAMBLE
\RequirePackage{color} %DIF PREAMBLE
\lstdefinelanguage{DIFcode}{ %DIF PREAMBLE
%DIF DIFCODE_UNDERLINE %DIF PREAMBLE
  moredelim=[il][\color{red}\sout]{\%DIF\ <\ }, %DIF PREAMBLE
  moredelim=[il][\color{blue}\uwave]{\%DIF\ >\ } %DIF PREAMBLE
} %DIF PREAMBLE
\lstdefinestyle{DIFverbatimstyle}{ %DIF PREAMBLE
	language=DIFcode, %DIF PREAMBLE
	basicstyle=\ttfamily, %DIF PREAMBLE
	columns=fullflexible, %DIF PREAMBLE
	keepspaces=true %DIF PREAMBLE
} %DIF PREAMBLE
\lstnewenvironment{DIFverbatim}[1][]{\lstset{style=DIFverbatimstyle}}{} %DIF PREAMBLE
\lstnewenvironment{DIFverbatim*}[1][]{\lstset{style=DIFverbatimstyle,showspaces=true}}{} %DIF PREAMBLE
\lstset{extendedchars=\true,inputencoding=utf8}

%DIF END PREAMBLE EXTENSION ADDED BY LATEXDIFF

\usepackage{soul}
\usepackage{xcolor}
\sethlcolor{yellow}
\providecommand{\DIFaddbegin}{}
\begin{document}

\section{Response to the Academic
Editor}\label{response-to-the-academic-editor}

\textbf{Manuscript:} \emph{Proxy-Based Diagnostics of Quantum Oracle
Sketching Robustness for Non-IID Sensor and Telemetry Streams}\\
\textbf{Journal:} \emph{Sensors} (MDPI) · \textbf{Manuscript ID:}
sensors-4470240

\textbf{Editor's comment.} \emph{``I strongly invite the Authors to not
include preprints among references. These works are still undergoing the
peer-review process, thus meaning that their content may change prior to
publication (if they will be finally published). The Authors can
substitute them with related works already published.''}

\textbf{Response.} We thank the Academic Editor \DIFdelbegin \DIFdel{and agree with the principle. We audited all }\DIFdelend \DIFaddbegin \DIFadd{for this helpful
recommendation. In accordance with the Editor's request, we have removed
the Zhao et al.~(2026) preprint, ``Exponential quantum advantage in
processing massive classical data'' (arXiv:2604.07639), from the revised
manuscript and from the reference list. The revised bibliography
contains }\DIFaddend 54 \DIFdelbegin \DIFdel{references: every entry was checked against
the publisher's own metadata, and all 51 DOIs resolved and matched the cited title, first author, venue, and year. The audit found two
preprints.

One had in fact been published
and has been replaced: the Kallaugher, }\DIFdelend \DIFaddbegin \DIFadd{references and no preprints.

}

\DIFadd{Before editing, we audited every in-text use of the preprint together
with the directly dependent clauses. We also examined the published
works already cited in the manuscript --- Kallaugher (FOCS 2021),
Kallaugher, }\DIFaddend Parekh and Voronova \DIFdelbegin \DIFdel{paper now carries its published version
(}\DIFdelend \DIFaddbegin \DIFadd{(STOC 2024 and }\DIFaddend SOSA 2025\DIFdelbegin \DIFdel{, pp.~9--45, DOI 10.1137/1.9781611978315.2). The same sweep
also fixed a transposed article number in the Su, Guo and Wang (2024)
entry, which now reads }\emph{\DIFdel{Information Sciences}} %DIFAUXCMD
\textbf{\DIFdel{676}}%DIFAUXCMD
\DIFdel{, article
120799, DOI 10.1016/j.ins.2024.120799.

}%DIFDELCMD < 

%DIFDELCMD < %%%
\DIFdel{One preprint remains: Zhao }\DIFdelend \DIFaddbegin \DIFadd{), and Huang }\DIFaddend et
al.\DIFdelbegin \DIFdel{, ``Exponential quantum advantage in
processing massive classical data'' (arXiv:2604.07639). On 16 August
2026 we re-checked whether it has been published or accepted anywhere.
It has not: the arXiv record is unchanged, no bibliographic record we
identified indicates a
published or in-press version, and the authors'
own publication pages still list only the arXiv posting. We retain it
because it is not cited as evidentiary support for a claim of ours; itis the specific framework under study. The paper examines how that
framework's correlation assumptions behave on non-IID streams; citing a
related published work in its place would attribute its theorems }\DIFdelend \DIFaddbegin \DIFadd{~(}\emph{\DIFadd{Science}} \DIFadd{2022). These establish important but different
results: problem-specific quantum-streaming space advantages, a
black-box quantum-sketch construction, and quantum advantage for
learning properties of quantum systems. None of them establishes the
framework-specific classification, dimension-reduction, correlated-data,
repetition-overhead, refreshing-time, or classical-hardness claims for
which the preprint had been cited. We therefore did not substitute any
of them for it, since doing so would attribute results }\DIFaddend to authors who
did not \DIFdelbegin \DIFdel{prove them, and removing it would leave the
framework the paper analyses unattributed}\DIFdelend \DIFaddbegin \DIFadd{obtain them}\DIFaddend .

\DIFdelbegin \DIFdel{We have, however, reduced the manuscript's reliance on it to the minimum
necessary for accurate attribution. Its in-text citations were cut from
eight to six by consolidating two repeats. The entry is labelled
``Preprint'' and is formatted in accordance with the MDPI Chicago
reference style for unpublished manuscripts and working papers, and a
footnote at its first citation states this status and explains why the preprint is not used as evidentiary support for the paper's empirical
conclusions: no quantum algorithm is implemented or simulated, and every
inherited quantity is listed with its provenance and epistemic status in Table 1. We also added a strong peer-reviewed published anchor of the same general kind --- }\DIFdelend \DIFaddbegin \DIFadd{Instead, only the affected statements were revised. Framework-specific
theorem and provenance claims were removed; each remaining reference is
now cited only for the result it actually establishes; and the
operational quantities τ, r, T\_eff, ε and the accuracy proxy are stated
unambiguously as constructs of the present paper. Section 3.3
accordingly separates published theory and its boundary, this paper's
assumptions, its heuristic constructions, its illustrative memory
references, and its scope limitations, and states explicitly that the
published quantum-streaming results do not establish the operational
model used here and that this paper adds no quantum theorem. Table 1's
provenance row and the associated memory-reference descriptions are
relabelled as paper-defined illustrative scaling references. The
remaining Zhao-dependent clauses in the abstract and introduction were
minimally corrected.

}

\DIFadd{The reference audit that accompanied this revision also confirmed the
published form of the }\DIFaddend Kallaugher, Parekh and Voronova \DIFdelbegin \DIFdel{(STOC 2024}\DIFdelend \DIFaddbegin \DIFadd{work (SOSA 2025}\DIFaddend ,
pp.~\DIFdelbegin \DIFdel{1805--1815, DOI }\DIFdelend \DIFaddbegin \DIFadd{9--45, DOI 10.1137/1.9781611978315.2), added the published STOC 2024
result (DOI }\DIFaddend 10.1145/3618260.3649709) \DIFdelbegin \DIFdel{, which establishes an
}\DIFdelend \DIFaddbegin \DIFadd{for the }\DIFaddend exponential
quantum--classical space separation \DIFdelbegin \DIFdel{for a natural streaming
problem }\DIFdelend --- \DIFdelbegin \DIFdel{co-cited with it in the Related Work and in the footnote. In
doing so, we }\DIFdelend \DIFaddbegin \DIFadd{which }\DIFaddend corrected a Related Work
sentence that \DIFdelbegin \DIFdel{had }\DIFdelend \DIFaddbegin \DIFadd{would otherwise have }\DIFaddend credited the 2021 FOCS paper with an
exponential separation\DIFdelbegin \DIFdel{; the FOCS 2021 result
provides a polynomial quantum advantage}\DIFdelend , whereas the \DIFdelbegin \DIFdel{STOC }\DIFdelend \DIFaddbegin \DIFadd{advantage established there is
polynomial --- and corrected a transposed article number in the Su, Guo
and Wang (}\DIFaddend 2024\DIFdelbegin \DIFdel{work
establishes the exponential separation. The revised bibliography now
contains 55 entries, including this newly added published anchor, and we
re-checked that every in-text citation matches exactly one bibliography
entry}\DIFdelend \DIFaddbegin \DIFadd{) entry, now }\emph{\DIFadd{Information Sciences}} \textbf{\DIFadd{676}}\DIFadd{,
article 120799, DOI 10.1016/j.ins.2024.120799}\DIFaddend .

\DIFdelbegin \DIFdel{Should the work appear in a peer-reviewed venue before this manuscript
reaches proofs, we will substitute the published version at that stage.
If the Editor considers the presence of even this single, explicitly
labelled preprint unacceptable, we would be grateful for the Editor's
preferred remedy and will follow it}\DIFdelend \DIFaddbegin \DIFadd{No replacement reference was added to preserve the citation count, and
no experimental data, methodology, equations, proxy definitions,
datasets, seeds, numerical results, figure artwork or conclusions were
changed. All reviewer-requested revisions remain in place, every in-text
citation was re-checked against the bibliography, and the manuscript
compiles with no unresolved citations or references}\DIFaddend .

On behalf of all authors,

\textbf{AbdelMoniem Helmy} (corresponding author)
\href{mailto:abdelmoniem.hafez@cu.edu.eg}{\nolinkurl{abdelmoniem.hafez@cu.edu.eg}}
· ORCID 0000-0001-5996-6019

\end{document}
